% authority: science-draft
% status: proposal-only
% route: ontology-law-research-packet
% trigger_classification: derivation_critical_missing_source_law
% adoption_status: blocked_adoption_open_continuation
\documentclass[11pt]{article}
\usepackage[margin=1in]{geometry}
\usepackage{amsmath,amssymb,amsthm}
\usepackage{booktabs}
\usepackage{enumitem}

\newtheorem{definition}{Definition}
\newtheorem{theorem}{Theorem}
\newtheorem{proposition}{Proposition}
\newtheorem{remark}{Remark}
\newcommand{\Ftwo}{\mathbb F_2}
\newcommand{\Sel}{\mathsf{CycleBoundarySel}_{\mathrm{src}}}
\newcommand{\AdmVar}{\mathsf{AdmVar}_{\mathrm{src}}}
\newcommand{\EqCand}{\mathsf{EqSrc}^{\mathrm{cand}}}
\newcommand{\Bottom}{\mathsf{Bottom}_{\mathrm{src}}}

\title{Proposal-Only EqSrc Cycle--Boundary Selector Law Candidate v1:\\
Marked-Source Naturality, Fixed-Carrier Discrimination, and Robustness}
\author{Ontology Formalizer control packet RT-20260718-035}
\date{2026-07-19}

\begin{document}
\maketitle

\section{Control Status}

This packet defines
\(\mathsf{EqSrcCycleBoundarySelectorLaw}^{\mathrm{cand},v1}_{\mathrm{src}}\)
as a proposal-only source-extension object. Current ontology does not derive
the unary source mark used below, does not select this marked-record category,
and does not establish the resulting translations as physically admissible.
The status is \texttt{proposal-only}, \texttt{source-extension data}, and
\texttt{blocked\_adoption\_open\_continuation}.

\begin{verbatim}
candidate_result: candidate_formalized_pending_fresh_smuggling_audit
candidate_name: EqSrcCycleBoundarySelectorLaw_src^cand,v1
formal_model: finite_odd_oriented_marked_source_records_over_F2
partial_selector: CycleBoundarySel_src(S)=(A_S,B_S) or tagged_bottom
cycle_space: A_S=(F2^X)^Aut(X,P)
boundary_space: B_S=ker(ell_P restricted to A_S)
candidate_eqsrc: kernel_pair_of_q_S:A_S->A_S/B_S
variation_definition_order: source_mark_then_u_then_variations_then_A_then_B_then_relation
naturality: marked_source_isomorphism_natural
automorphism_equivariance: proved
relation_level_uniqueness: unique_up_to_quotient_codomain_bijection
mark_and_parity_selector_uniqueness: not_claimed
chain_presentation_uniqueness: not_claimed
finite_witness: four_point_singleton_mark
witness_state_count: 4
witness_class_count: 2
witness_class_size: 2
fixed_carrier_alternate_mark_relation_changed: true
current_ontology_derives_candidate: false
physical_admissibility_established: false
candidate_status: proposal-only
adoption_status: blocked_adoption_open_continuation
general_EqSrc_discharged: false
distance_to_gr_ledger_changed: false
freeze_decision: not_frozen
next_route: fresh_smuggling_auditor_review
\end{verbatim}

The exact
\(\mathsf{EqSrcIntrinsicDiscriminatorAdmissibilityLaw}^{\mathrm{cand},v3}_{\mathrm{src}}\)
construction--audit--stress cycle remains locally frozen. The present object
moves upstream: it makes the missing source mark explicit and selects only the
quotient-relevant pair before any candidate relation is evaluated.

\section{Marked Source-Record Category and Partial Selector}

\begin{definition}[Odd-oriented-marked source record]
An admitted record is a pair \(S=(X,P)\), where \(X\) is a finite set and
\[
 \varnothing\neq P\subsetneq X,\qquad |P|\equiv1\pmod2.
\]
A morphism \(f:S\to S'\) is a bijection \(f:X\to X'\) satisfying
\(f(P)=P'\). These objects and morphisms form the groupoid
\(\mathsf{FinMark}^{\mathrm{odd},*}_{\mathrm{src}}\).

The oriented mark \(P\), rather than the unordered partition
\(\{P,X\setminus P\}\), is displayed source data. It is not inferred from a
desired relation, target geometry, benchmark behavior, or workflow state.
\end{definition}

\begin{definition}[Pre-relation variation datum]
For \(S=(X,P)\), set
\[
 V_S=\Ftwo^X,\qquad u_S=\mathbf1_{X\setminus P}.
\]
Before any fixed space, boundary space, quotient, or relation is evaluated,
define the affine involutions
\[
 J_S(\epsilon):V_S\longrightarrow V_S,\qquad
 J_S(\epsilon)(a)=a+\epsilon u_S,\qquad \epsilon\in\Ftwo.
\]
This gives an affine \(\Ftwo\)-action because
\(J_S(\epsilon)J_S(\delta)=J_S(\epsilon+\delta)\).
\end{definition}

\begin{definition}[Cycle--boundary selector]
Set
\[
 G_S=\operatorname{Aut}(X,P),\qquad A_S=V_S^{G_S}.
\]
Write
\[
 \ell_{P,S}:V_S\longrightarrow\Ftwo,\qquad
 \ell_{P,S}(a)=\sum_{x\in P}a_x,
\]
and define
\[
 B_S=\ker\!\left(\ell_{P,S}|_{A_S}\right).
\]
The partial law returns
\[
 \Sel(S)=(A_S,B_S\subsetneq A_S)
\]
when \(S\) is admitted and \(0\subsetneq B_S\subsetneq A_S\); otherwise it
returns the tagged value \(\Bottom(S)\) and defines no candidate relation.
\end{definition}

\begin{proposition}[Exact form on the admitted domain]
Let \(\mathbf 1_P,\mathbf 1_{X\setminus P}\in V_S\) be the two orbit
indicators. Then
\[
 A_S=\operatorname{span}_{\Ftwo}
 \{\mathbf1_P,\mathbf1_{X\setminus P}\},\qquad
 B_S=\operatorname{span}_{\Ftwo}\{\mathbf1_{X\setminus P}\}.
\]
Consequently \(\dim A_S=2\), \(\dim B_S=1\), and every admitted record
passes the nondegeneracy test.
\end{proposition}

\begin{proof}
The automorphism group is the product of the full symmetric groups on \(P\)
and \(X\setminus P\), so a fixed vector is constant on each of those two
orbits. Since \(|P|\) is odd,
\(\ell_{P,S}(\mathbf1_P)=1\) and
\(\ell_{P,S}(\mathbf1_{X\setminus P})=0\). The displayed kernel follows.
\end{proof}

\section{Dependency Order and Independent Variations}

\begin{definition}[Admitted variation class]
Restrict the already-defined affine action to \(A_S\):
\[
 \AdmVar(S)=\{J_S(\epsilon)|_{A_S}:\epsilon\in\Ftwo\}.
\]
The exact-form proposition proves \(u_S\in A_S\) and
\(B_S=\Ftwo u_S\), so these restrictions are affine automorphisms of
\(A_S\). They are not asserted to be linear endomorphisms.
\end{definition}

The dependency order is
\[
 (X,P)\longrightarrow(V_S,u_S,J_S)
 \longrightarrow(G_S,A_S,\ell_{P,S})\longrightarrow B_S
 \longrightarrow q_S\longrightarrow\EqCand_S.
\]
Neither \(\EqCand_S\), its stabilizer, a target object, benchmark behavior,
nor process authority occurs upstream of the affine action, \(B_S\), or the
candidate relation. This is a formal non-circularity certificate relative to
the proposed marked record.
It does not derive the mark from current ontology and does not establish
physical admissibility.

\section{Candidate Relation and Relation-Level Uniqueness}

\begin{definition}[Candidate relation]
Let
\[
 q_S:A_S\longrightarrow A_S/B_S
\]
be the quotient map. Define
\[
 a\,\EqCand_S\,a'
 \quad\Longleftrightarrow\quad
 q_S(a)=q_S(a')
 \quad\Longleftrightarrow\quad a-a'\in B_S.
\]
\end{definition}

\begin{theorem}[Relation-level quotient uniqueness]
Let \(d:A_S\to D\) be a surjection such that
\[
 d(a)=d(a')\quad\Longleftrightarrow\quad a-a'\in B_S.
\]
There is a unique bijection \(\bar d:A_S/B_S\to D\) with
\(d=\bar d\circ q_S\). Thus the selected relation is unique up to a unique
relabeling of the quotient codomain under the stated same-kernel hypothesis.
\end{theorem}

\begin{proof}
Set \(\bar d([a])=d(a)\). Equality modulo \(B_S\) makes this well defined;
the exact kernel-pair hypothesis makes it injective; surjectivity is inherited
from \(d\). Any factorization must take \([a]\) to \(d(a)\).
\end{proof}

This is not uniqueness of the mark, the source record, or every possible
selector. In particular, source symmetry alone does not select the oriented
mark or the parity functional \(\ell_{P,S}\). The theorem is uniqueness of a
surjective label for one already selected kernel pair.

\section{Naturality and Automorphism Equivariance}

\begin{theorem}[Marked-source naturality]
For a marked bijection \(f:S\to S'\), let
\[
 F_f:V_S\longrightarrow V_{S'},\qquad
 (F_fa)_{f(x)}=a_x.
\]
Then
\[
 F_f(A_S)=A_{S'},\qquad F_f(B_S)=B_{S'},
\]
and the induced quotient bijection \(\overline F_f\) makes
\[
 \boxed{\quad q_{S'}F_f=\overline F_fq_S\quad}
\]
commute on \(A_S\). Consequently marked-source isomorphisms preserve and
reflect \(\EqCand\), and composition, identity, and inverse certificates hold.
\end{theorem}

\begin{proof}
Conjugation by \(f\) identifies \(G_S\) with \(G_{S'}\), so \(F_f\) maps
fixed vectors bijectively. Because \(f(P)=P'\),
\(\ell_{P',S'}F_f=\ell_{P,S}\), hence \(F_f\) maps the restricted kernels
bijectively. The quotient square and functorial certificates follow.
\end{proof}

For \(f\in G_S\), every \(a\in A_S\) is fixed by definition, so the
automorphism-equivariance claim is exact but weak: the selected state space
is the fixed space, so the restricted action is pointwise trivial. This
naturality is only for the displayed marked-source groupoid. It says nothing
about unmarked carrier bijections, mark-changing maps, or physical covariance
under an independently justified category of source changes.

\section{Chain-Presentation Compatibility}

The selected pair has a canonical minimal two-truncated presentation
\[
 C_2(S)=B_S\xrightarrow{\iota_S}C_1(S)=A_S
 \xrightarrow{0}C_0(S)=0.
\]
Its cycle space is \(A_S\) and its boundary image is \(B_S\). More generally,
any typed presentation whose cycle space is \(A_S\) and whose boundary image
is \(B_S\) induces the same quotient and kernel pair.

No uniqueness of \(C_2\), a differential, or a full chain presentation is
claimed. The selected pair is authoritative only inside this proposal; a
differential is a presentation of \(B_S\), not a physically selected input.
An externally supplied presentation with a different boundary image is tagged
\texttt{presentation\_incompatible}; it does not make the canonical partial
selector return \(\Bottom\).

\section{Conditional Robustness}

\begin{theorem}[Finite admitted-variation robustness]
For every \(\epsilon\in\Ftwo\) and \(a\in A_S\),
\[
 q_S(J_S(\epsilon)a)=q_S(a).
\]
Every finite composition of admitted translations therefore acts trivially
on \(A_S/B_S\) and preserves \(\EqCand_S\).
\end{theorem}

\begin{proof}
The difference \(J_S(\epsilon)a-a=\epsilon u_S\) lies in
\(B_S=\Ftwo u_S\). Closure under finite composition follows from the affine
action law.
\end{proof}

This theorem is intentionally conditional. It is algebraic robustness for a
variation class derived from the displayed mark before the relation is
evaluated. It neither proves that nature supplies the mark nor that all
physically relevant variations are admitted.

\section{Four-Point Witness and Fixed-Carrier Discrimination}

Let \(X=\{0,1,2,3\}\). For \(S_0=(X,\{0\})\), write vectors in coordinate
order \(0,1,2,3\). Then
\[
 A_{S_0}=\{0000,1000,0111,1111\},\qquad
 B_{S_0}=\{0000,0111\}.
\]
The four selected source states form two two-element classes:
\[
\begin{array}{c|c}
\text{quotient label}&\text{class}\\\hline
0&\{0000,0111\}\\
1&\{1000,1111\}.
\end{array}
\]
The nonidentity admitted translation toggles \(0111\) and leaves the quotient
label fixed.

On the same carrier, use the distinct record
\(S_1=(X,\{1,2,3\})\). It has the same invariant state space but
\[
 B_{S_1}=\{0000,1000\},
\]
with classes \(\{0000,1000\}\) and \(\{0111,1111\}\). Therefore the two
explicit source marks select different boundary subspaces and different
relations on the same carrier. This is a structural marked-record mutation,
not robustness under every source variation and not evidence that current
ontology selects either mark. If the physical source gauge may exchange or
forget the two orientations while preserving the carrier, the candidate
fails closed until extra source-side provenance distinguishes them.

\section{Fail-Closed Branches}

The proposed law returns \(\Bottom\) and defines no candidate relation when:
\begin{enumerate}[leftmargin=*]
  \item the unary mark is absent, empty, full, or even-cardinality;
  \item the mark is inferred from a desired relation, its stabilizer, target
  geometry, benchmark success, empirical detector semantics, or process data;
  \item an alleged morphism is not a bijection or does not preserve the mark;
  \item the fixed-space or restricted-kernel calculation does not yield
  \(0\subsetneq B_S\subsetneq A_S\);
  \item an admitted variation is defined after or from the selected relation;
  \item naturality, composition, inverse, quotient-factorization, or
  relation-level uniqueness certification fails;
  \item an externally supplied chain presentation has the wrong cycle space
  or boundary image, in which case it is tagged
  \texttt{presentation\_incompatible} without changing the canonical
  selector output; or
  \item any generated derivative, registry, role, validator, handoff,
  checkpoint, or file ordering is treated as scientific authority.
\end{enumerate}
Failure closes only that candidate instance. It does not prove a general
EqSrc no-go or future source-extension impossibility.

\section{No-Target and No-Process Guard}

The construction depends only on a finite set, an explicit unary source mark,
finite permutations, \(\Ftwo\) linear algebra, and quotient formation.
Target atlas, topology, metric, proper time, detector semantics, desired
exact-GR behavior, registry metadata, role identity, validator output,
handoff text, checkpoint status, and generated derivatives are excluded from
the definition tree.

The mark is nevertheless a new underived source primitive candidate. Calling
the construction source-only does not make that mark ontology-derived or
physically admissible.

\section{Distance-to-GR, Freeze, and Next Route}

The result is
\texttt{candidate\_formalized\_pending\_fresh\_smuggling\_audit}. It supplies
a partial selector definition, naturality and uniqueness theorems, a
presentation lemma, conditional robustness, and a fixed-carrier witness. It
does not discharge general EqSrc and does not change the Distance-to-GR or
metric-use ledgers.

The exact v3 construction--audit--stress route remains locally frozen. This
distinct marked-source selector route is \texttt{not\_frozen} because it makes
the missing primitive explicit and provides a new theorem and witness payload.

The next lawful packet is one fresh bounded
\texttt{smuggling-auditor@0.2.0} audit of this exact candidate. It must test
whether the mark merely renames the missing differential choice, whether the
record category or variation class imports target or process data, whether
formal naturality is overread as physical covariance, and whether the finite
witness is overgeneralized. It may not repair or adopt the candidate.

\section{Forbidden Conclusions}

This artifact does not authorize canonical ontology modification, source-law
adoption, derivation of the unary mark, physical-admissibility establishment,
general EqSrc discharge, RetainH or GenH adoption, \(M_{\mathrm{src}}\)
adoption, \(g_{\mathrm{eff}}\) scope expansion, matter coupling, Einstein
equations, benchmark promotion, Gate Chair verdict, completed derivation,
future source-extension impossibility, program-level no-go, or global theory
rejection.

\section{Source Materials}

\begin{itemize}[leftmargin=*]
  \item Handoff 0759 and the RT-20260718-034 selector artifact.
  \item The registered Ontology Formalizer role contract v0.2.0.
  \item The source-extension classification checklist and no-target-import
  guard map.
  \item Finite group-action and quotient arguments reproduced above.
\end{itemize}

\end{document}
